\section{Introducción}
    Entre las operaciones de preprocesamiento para los problemas 
    de localización y segmentación se encuentra la detección de 
    bordes, esta técnica permite aislar de forma preliminar los 
    objetos que están presentes en una imagen. Por lo tanto, el 
    problema de los bordes no solamente se limita a una curiosidad 
    en visión computacional, sino que se vuelve un proceso intermedio 
    para mejorar y optimizar los modelos de IA más sofisticados 
    como las CNN (Convolutional Neural Network). Esta se vuelve la 
    principal motivación para entender cómo surgen los bordes, 
    cómo se pueden detectar y, aún más importante, cómo se pueden 
    integrar estas ideas en redes neuronales.\\

    Para realizar este proyecto, se hace uso de los filtros de Prewitt, 
    Sobel, Laplace, Laplaciano de Gaussiana y Canny para compararlos 
    en diferentes escenarios de bordes e imágenes con y sin suavizado 
    gaussiano con el fin de entender sus ventajas, desventajas y 
    cómo los bordes son afectados (su definición y ruido presente) por 
    el filtro específico que se emplea. \\

    De forma general, los filtros de Prewitt y Sobel permiten detectar 
    bordes geométricos sin dificultad, mientras que los filtros de Laplace 
    y Laplaciano de Gaussiana les va mejor en bordes curvos o no geométricos \parencite{tinoco2026pdi}. 
    El filtro de Canny emplea de los filtros de Prewitt o de Sobel para 
    detectar bordes horizontales, verticales y diagonales, esto favorece a 
    que este filtro sea el más conveniente para todos los escenarios, en 
    teoría \parencite{wiki2026canny}.